% =====================================================================
\chapter{Regla de la cadena y jacobiano}
\label{tema:2}
% =====================================================================

Una red neuronal es una composición de muchas funciones sencillas, y su pérdida depende de cada parámetro a través de largas cadenas de cálculos intermedios. Este tema generaliza la derivada a funciones con valores vectoriales mediante el jacobiano, demuestra la regla de la cadena en forma matricial y muestra cómo se organiza el cálculo sobre un grafo computacional, en modo directo o en modo inverso. Se sigue~\cite[§5.2, §5.3 y §5.6]{mml}, la clase~10 de~\cite{mit1802} y~\cite[§6.5]{gbc}.

\section{Funciones vectoriales y jacobiano}

\begin{definicion}[Jacobiano]\label{def:jacobiano}
Sea $U\subseteq\R^{n}$ abierto y $g=(g_1,\dots,g_m)\T\colon U\to\R^{m}$. La función $g$ es \emph{diferenciable} en $x$ si existe una matriz $J\in\R^{m\times n}$ tal que
\[
g(x+h)=g(x)+J\,h+o(\norm{h})\qquad\text{cuando } h\to 0.
\]
La matriz $J$ se llama \emph{jacobiano} de $g$ en $x$ y se escribe $J_g(x)$. Aplicando la \cref{prop:grad-parciales} a cada componente, $g$ es diferenciable si y solo si lo es cada $g_j$, y
\[
\bigl(J_g(x)\bigr)_{ji}=\frac{\partial g_j}{\partial x_i}(x),\qquad j=1,\dots,m,\quad i=1,\dots,n.
\]
La fila $j$ de $J_g(x)$ es $\nabla g_j(x)\T$; la columna $i$ es la derivada de $g$ respecto de $x_i$.
\end{definicion}

El jacobiano tiene una fila por cada salida y una columna por cada entrada. Para $m=1$ es una matriz fila y $\nabla f(x)=J_f(x)\T$.

\addtocounter{definicion}{1}%

\addtocounter{definicion}{1}%

\section{La regla de la cadena}

\begin{teorema}[Regla de la cadena]\label{teo:cadena}
Sean $g\colon U\subseteq\R^{n}\to\R^{m}$ diferenciable en $x$ y $f\colon V\subseteq\R^{m}\to\R^{k}$ diferenciable en $g(x)$, con $g(U)\subseteq V$. Entonces $f\circ g$ es diferenciable en $x$ y
\[
J_{f\circ g}(x)=J_f\bigl(g(x)\bigr)\,J_g(x).
\]
\end{teorema}

\begin{corolario}[Caso escalar]\label{cor:cadena-escalar}
Si $f\colon\R^{m}\to\R$ y $u=g(x)$, entonces
\[
\frac{\partial (f\circ g)}{\partial x_i}(x)=\sum_{j=1}^{m}\frac{\partial f}{\partial u_j}(u)\,\frac{\partial u_j}{\partial x_i}(x),
\qquad\text{es decir,}\qquad
\nabla_x (f\circ g)=J_g(x)\T\,\nabla_u f(u).
\]
\end{corolario}

La suma del \cref{cor:cadena-escalar} tiene una lectura directa: si $x_i$ influye en $f$ a través de varias variables intermedias $u_j$, su efecto total es la suma de los efectos transmitidos por cada una de ellas.

\addtocounter{definicion}{1}%

\section{Grafos computacionales}

\begin{definicion}[Grafo computacional]\label{def:grafo}
Un \emph{grafo computacional} es un grafo dirigido acíclico cuyos nodos son variables: los nodos sin predecesores son las entradas, y cada uno de los demás nodos $v$ se calcula como una función diferenciable $v=\phi_v(\text{padres}(v))$ de sus predecesores. Si la salida es un escalar $f$, se llama \emph{adjunto} de un nodo $v$ a $\bar v=\partial f/\partial v$, la derivada de $f$ respecto de $v$ considerando todas las variables posteriores como funciones de $v$~\cite[§6.5.1]{gbc}.
\end{definicion}

\begin{proposicion}[Retropropagación]\label{prop:retro}
En un grafo computacional con salida escalar $f$ se cumple $\bar f=1$ y, para todo nodo $v\neq f$,
\[
\bar v=\sum_{w\,:\,v\in\text{padres}(w)}\bar w\;\frac{\partial\phi_w}{\partial v}.
\]
Por tanto, recorriendo los nodos en orden topológico inverso se obtienen los adjuntos de todas las entradas, es decir, el gradiente de $f$.
\end{proposicion}

\addtocounter{figure}{1}%

\section{Modo directo y modo inverso}

Sea una composición $F=f_L\circ\dots\circ f_1$, con $f_\ell\colon\R^{d_{\ell-1}}\to\R^{d_\ell}$. Por el \cref{teo:cadena},
\[
J_F=J_L\,J_{L-1}\cdots J_1,\qquad J_\ell\in\R^{d_\ell\times d_{\ell-1}},
\]
con cada $J_\ell$ evaluado en el valor intermedio correspondiente. El producto de matrices es asociativo, pero el coste de calcularlo depende del orden en que se agrupen los factores~\cite[§5.6]{mml}.

\begin{proposicion}[Coste de los dos modos]\label{prop:modos}
Multiplicar una matriz $a\times b$ por otra $b\times c$ requiere $abc$ productos escalares. En consecuencia:
\begin{enumerate}
  \item en \emph{modo directo}, se calcula $J_F v=J_L(\cdots(J_2(J_1 v)))$ para un vector $v\in\R^{d_0}$ con $\sum_\ell d_\ell d_{\ell-1}$ productos; obtener el jacobiano completo exige repetirlo $d_0$ veces, una por columna;
  \item en \emph{modo inverso}, se calcula $w\T J_F=((w\T J_L)J_{L-1})\cdots J_1$ para un vector $w\in\R^{d_L}$ con el mismo número de productos; obtener el jacobiano completo exige repetirlo $d_L$ veces, una por fila.
\end{enumerate}
Si la salida es escalar ($d_L=1$), el gradiente se obtiene con una sola pasada en modo inverso.
\end{proposicion}

La retropropagación de la \cref{prop:retro} es el modo inverso aplicado a un grafo. Su coste es un múltiplo pequeño del coste de evaluar $f$, independiente del número de entradas; por ello la pérdida de un modelo con miles de millones de parámetros puede derivarse con una sola pasada hacia atrás. El precio es la memoria: el paso hacia atrás necesita los valores intermedios del paso hacia delante. La implementación de este mecanismo (diferenciación automática) es el objeto del módulo~10.

\addtocounter{definicion}{1}%

